% !TeX spellcheck = ru_RU
% !TEX root = vkr.tex


\section*{Заключение}

В ходе выполнения выпускной квалификационной работы получены следующие результаты.

\begin{enumerate}
	\item Выполнен обзор существующих методов обнаружения цифровых сигналов и особенностей их реализации на программируемых логических интегральных схемах. Установлено, что корреляционный метод является теоретически обоснованным и практически применимым для обнаружения сигналов с известной структурой преамбулы. Для аппаратной реализации выбран подход, основанный на вычислении взаимнокорреляционной функции с использованием быстрого преобразования Фурье в соответствии с теоремой о свёртке, что обеспечивает эффективность реализации.
	
	\item Спроектирована архитектура ПЛИС-подсистемы обнаружения, предназначенной для интеграции в устройство Pluto SDR+. Разработанная архитектура реализует вычисление взаимнокорреляционной функции входного сигнала с эталонной преамбулой для набора частотных гипотез и учитывает особенности существующего цифрового тракта приёма и передачи данных между логической частью (PL) и процессорной системой (PS).
	
	\item Выполнена реализация обнаружителя на языке VHDL для ПЛИС, входящей в состав SoC семейства Xilinx Zynq-7000. Разработанный модуль интегрирован в штатную HDL-прошивку устройства Pluto SDR+. Реализация выполнена с использованием потоковой обработки данных и интерфейса AXI4-Stream, обеспечивающего передачу результатов корреляционной обработки в тракт DMA.
	
	\item Проведено функциональное тестирование разработанной подсистемы средствами RTL-моделирования в среде Vivado. По результатам симуляции подтверждена корректность вычисления взаимнокоррреляционной функции и работоспособность реализованного решения. Полученные данные сопоставлены с результатами программного моделирования, что показало их согласованность.
\end{enumerate}





%\section*{Заключение}
%
%В ходе работы были достигнуты следующие результаты.
%\begin{enumerate}
%\item Выполнен обзор литературы, посвященной обнаружению сигналов. Установлено, что наиболее распространенным является метод корреляционного обнаружения сигналов. Сведений о реализации на ПЛИС алгоритмов для этого метода найдено не было. Выполнен расчет вычислительной сложности базового корреляционного алгоритма; показано, что непосредственная реализация данного алгоритма в ПЛИС будет обеспечивать скорость работы, не представляющую практического интереса. Выполнено экспериментальное исследование базового корреляционного алгоритма и алгоритма НИИРТ; показано, что при фиксированном диапазоне частотных ошибок алгоритмы демонстрируют сопоставимое качество обнаружения, однако реализации на ПЛИС алгоритма НИИРТ требует больше ресурсов. Для реализации на ПЛИС выбран базовый корреляционный алгоритм.
%\item Выполнена модификация базового корреляционного алгоритма, заключающаяся в применении операции циклического сдвига к спектру входного сигнала для определения величины частотной ошибки входного сигнала.
%\item Спроектирована архитектура обнаружителя, включающая буфер входного сигнала, хранилище преамбул, блок расчета корреляции и пороговое устройство.
%\item Выполнена реализация обнаружителя на языке VHDL для модели ПЛИС из состава SoC Xilinx Zynq-7000, включая решение следующих задач:
%    \begin{enumerate}
%    \item синхронизация входного сигнала и блоков быстрого преобразования Фурье, тактируемых различными частотами;
%    \item выбор размеров буферов для отсчетов входного сигнала и для блоков быстрого преобразования Фурье в условиях ограниченности ресурсов памяти, с учетом требований к длине преамбулы и предельной частоте дискретизации входного сигнала;
%    \item реализация механизма частотного сдвига сигнала в условиях нехватки DSP элементов;
%    \item расчет непрерывной взаимнокорреляционной функции при наличии буферизации входного сигнала.
%    \end{enumerate}
%\item Средствами Vivado IDE было выполнено симуляционное тестирование на уровне RTL, которое показало корректность функционирования разработанного IP-ядра обнаружителя, продемонстрировав его способность достоверно обнаруживать тестовые сигналы с преамбулой в условиях заданного частотного рассогласования и помех. 
%\end{enumerate}








% В ходе работы за весенний семестр были достигнуты следующие результаты.
% \begin{enumerate}
%     \item Реализована математическая модель алгоритма обнаружения цифровых сигналов на языке Python.
%     \item На основе результатов моделирования осуществлен сбор предельных характеристик обнаружителя, при которых выполняется успешное обнаружение сигналов.
%     \item На основе собранных эмпирические данных осуществлен вывод математических выражений, являющихся аппроксимацией параметров алгоритма.
%     \item Осуществлена сборка и синтез прошивки ПЛИС на базе драйверов фирмы Analog Devices, обеспечивающую интерфейс управления микросхемой ad9361, а также интерфейс получения отсчетов оцифрованного сигнала.
%     \item На базе IP-ядер фирмы Xilinx реализовано IP-ядро переноса по частоте.
%     \item На базе TCL-скриптов реализована возможность получения оцифрованных данных от АЦП на ПК.
% \end{enumerate}


% Также сюда можно написать планы развития работы в будущем, или, если их много, выделить под это отдельную предпоследнюю главу.

